\scanpage{21}
anneau de valuation discrète, et où on se
donne une section \inserted{rationnelle} de $M$ sur $\unsure{B|S}$ i.e. que
on regarde l'unique $A_K$ abélien
sous-jacent à $\mathrm{Pic}^0(X_K)$\deleted{\unsure{,|$\ast$}} \unsure{;|?|$+$} prouver \unsure{qu'en|que}
la section est partout définie. Or par le
\uuline{théorème de Koizumi} $A_K$ ``se réduit bien''
i.e. provient d'un schéma abélien $A$
sur $V$. On va essayer de prolonger
en un \deleted{\unsure{hom}} morphisme $A\to\uline{\mathrm{Pic}}^{-}_{X/S}$
i.e. $A\to P$. C'est possible, grâce au
fait que les fibres de $P$ sur $S$ satisfont
les hypothèses du théorème de \unsure{connexion|connexions} \inserted{\unsure{linéaire|birationnelle}}
de Murre-Chow. [On n'a pas besoin de
ce théorème si $X$ simple sur $S$].
Reste la question si $u:A\to P\deleted{\unsure{x}}$ est une
\uline{immersion} fermée, d'image $P$. L'image a
\deleted{\unsure{est}} évidemment \inserted{\unsure{une fibre|ses fibres}} de dim égale à $\dim A_K$ \deleted{\unsure{génér}}
\inserted{$=\dim\mathrm{Pic}(X_K)$} \deleted{et} [\unsure{continuité} de la dimension des fibres
de $\mathrm{Ker}\,u$] donc de dim égale à
$\dim\mathrm{Pic}(X_k)$ [\uuline{théorème de l'invariance
de la dimension des fibres des Picards}]
% slanted
\margin{Attention, cela est-il vrai si $f:X\to S$ n'est pas simple ??}
% archive sheet number: 20
